\documentclass[a4paper,12pt]{article}
\usepackage[utf8]{inputenc}
\usepackage[italian]{babel}
\usepackage{listings}
\usepackage{xcolor}
\usepackage{float}
\usepackage{placeins}

\lstset{
  language=C,
  backgroundcolor=\color{lightgray!20},
  frame=single,
  breaklines=true,
  numbers=left,
  numberstyle=\tiny,
  stepnumber=1,
  numbersep=5pt,
  captionpos=b,
  basicstyle=\ttfamily\small,
  keywordstyle=\color{blue}\bfseries,
  commentstyle=\color{green!60!black},
  stringstyle=\color{red},
  showstringspaces=false,
  float=H
}

\title{LabReSiD26 - Hands-On 4}
\author{Davide Ferrara - 518629}
\date{}

\begin{document}

\maketitle

% ---------------------------------------------------------------------------
\section{Esercizio 1: Stampa leggibile del protocollo}
% ---------------------------------------------------------------------------

Il campo \texttt{ip->protocol} contiene il numero di protocollo del livello di trasporto
incapsulato nel pacchetto IP. Nel codice di partenza veniva stampato come numero intero
con \texttt{\%d}. La funzione \texttt{proto\_name()} converte il numero in una stringa
leggibile tramite uno switch case che matcha il numero del protocollo e ritorna la stringa
corrispondente:

\begin{lstlisting}[caption={proto\_name() -- implementazione}]
const char *proto_name(uint8_t protocol) {
  switch (protocol) {
  case 1:
    return "ICMP";
  case 6:
    return "TCP";
  case 17:
    return "UDP";
  default:
    return "sconosciuto";
  }
}
\end{lstlisting}

La \texttt{printf} è stata modificata per usare \texttt{proto\_name()} al posto di
\texttt{\%d}, mantenendo il numero tra parentesi per completezza:

\begin{lstlisting}[caption={Uso nel loop di stampa}]
printf("  PROTO: %s (%d)\n", proto_name(ip->protocol), ip->protocol);
\end{lstlisting}

\begin{lstlisting}[language={}, caption={Output}]
  PROTO: ICMP (1)
\end{lstlisting}

% ---------------------------------------------------------------------------
\section{Esercizio 2: Decodifica completa dell'header IP}
% ---------------------------------------------------------------------------

Ho esteso il programma per stampare tutti i campi significativi dell'header IP.
È importante segnalare che se \texttt{ip->ihl < 5} il pacchetto viene scartato
con \texttt{continue} nel ciclo, in quanto un valore inferiore a 5 indica un
header malformato (la lunghezza minima valida è $5 \times 4 = 20$ byte).

Il campo \texttt{frag\_off} è un \texttt{uint16\_t} che comprime tre informazioni:

\begin{center}
\texttt{| bit 15 (riservato) | bit 14 (DF) | bit 13 (MF) | bit 12--0 (offset) |}
\end{center}

\begin{itemize}
  \item \textbf{DF} (Don't Fragment, bit 14): il pacchetto non deve essere frammentato.
  \item \textbf{MF} (More Fragments, bit 13): esistono altri frammenti successivi.
  \item \textbf{Fragment Offset} (bit 12--0, 13 bit): posizione del frammento in unità da 8 byte.
\end{itemize}

Si estraggono con operazioni bitwise dopo aver convertito \texttt{frag\_off} in host byte
order con \texttt{ntohs()}:

\begin{lstlisting}[caption={Estrazione flag e offset da frag\_off}]
uint16_t fo = ntohs(ip->frag_off);
int flag_df = (fo >> 14) & 1; /* sposta bit 14 in pos. 0, prendi solo quello */
int flag_mf = (fo >> 13) & 1; /* stesso per bit 13 */
int offset  =  fo & 0x1FFF;   /* maschera i 13 bit bassi */
\end{lstlisting}

\begin{lstlisting}[caption={Decodifica completa header IP}]
if (ip->ihl < 5) {
  printf("[SNIFFER.c] Pacchetto malformato\n");
  continue;
}

uint16_t tot_len = ntohs(ip->tot_len);
int len = ip->ihl * 4;
uint16_t fo = ntohs(ip->frag_off);
int flag_df = (fo >> 14) & 1; /* Don't Fragment */
int flag_mf = (fo >> 13) & 1; /* More Fragments */
int offset  =  fo & 0x1FFF;   /* Fragment Offset (in unita' da 8 byte) */

printf("  VERSION: %d\n",    ip->version);
printf("  IHL: %d byte\n",   len);
printf("  LUNGHEZZA: %d byte\n", tot_len);
printf("  TTL: %d\n",        ip->ttl);
printf("  PROTO: %s (%d)\n", proto_name(ip->protocol), ip->protocol);
printf("  SRC: %s\n",        inet_ntoa(s));
printf("  DST: %s\n",        inet_ntoa(d));
printf("  FLAGS: %s%s\n",    flag_df ? "DF " : "", flag_mf ? "MF" : "");
printf("  FRAG OFF: %d\n\n", fo);
\end{lstlisting}

\begin{lstlisting}[language={}, caption={Output -- ping -c 3 8.8.8.8}]
sudo ./sniffer
In ascolto (solo ICMP)... premi Ctrl+C per uscire

Pacchetto #1  (84 byte)
  VERSION: 4
  IHL: 20 byte
  LUNGHEZZA: 84 byte
  TTL: 116
  PROTO: ICMP (1)
  SRC: 8.8.8.8
  DST: 192.168.8.230
  FLAGS:
  FRAG OFF: 0
\end{lstlisting}

\texttt{FLAGS} è vuoto e \texttt{FRAG OFF: 0} perché il pacchetto non è frammentato
e il flag DF non è settato nella risposta di \texttt{8.8.8.8}.

% ---------------------------------------------------------------------------
\section{Challenge: Naviga nel payload -- decodifica l'header ICMP}
% ---------------------------------------------------------------------------

La funzione \texttt{ip\_payload()} calcola il puntatore all'inizio del payload IP
tenendo conto delle eventuali opzioni (\texttt{ihl > 5}) ed esegue due controlli
di sicurezza prima di restituire il puntatore:

\begin{itemize}
  \item \texttt{ip->ihl < 5}: header malformato, lunghezza impossibile.
  \item \texttt{n < hlen}: pacchetto troncato — \texttt{recvfrom()} ha consegnato
        meno byte di quanti l'header dichiara. Su IP non esiste ritrasmissione,
        il pacchetto va scartato.
\end{itemize}

\begin{lstlisting}[caption={ip\_payload() -- implementazione}]
unsigned char *ip_payload(const unsigned char *buf, ssize_t n,
                          size_t *payload_len) {
  struct iphdr *ip = (struct iphdr *)buf;
  if (ip->ihl < 5)
    return NULL; /* malformato */
  int hlen = ip->ihl * 4;
  if (n < hlen)  /* se n < hlen e accedo a buf+hlen vado in segfault */
    return NULL; /* troncato: in IP non esiste ritrasmissione, lo scarto */
  *payload_len = n - hlen;
  return (unsigned char *)buf + hlen;
}
\end{lstlisting}

\begin{lstlisting}[caption={Uso nel loop -- decodifica header ICMP}]
size_t payload_len;
unsigned char *payload = ip_payload(buf, n, &payload_len);
if (payload == NULL)
  continue;

printf("  PAYLOAD:\n");
if (ip->protocol == 1) {
  struct icmphdr *icmp = (struct icmphdr *)payload;
  printf("  TYPE:     %d\n",  icmp->type);
  printf("  CODE:     %d\n",  icmp->code);
  printf("  CHECKSUM: %d\n",  ntohs(icmp->checksum));
  printf("  ID:       %d\n",  ntohs(icmp->un.echo.id));
  printf("  SEQ:      %d\n\n", ntohs(icmp->un.echo.sequence));
}
\end{lstlisting}

\begin{lstlisting}[language={}, caption={Output -- ICMP Echo Reply da 8.8.8.8}]
  PAYLOAD:
  TYPE: 0
  ID: 1423
  SEQ: 256
\end{lstlisting}

\texttt{TYPE: 0} corrisponde a Echo Reply — la risposta al ping. Per vedere
\texttt{TYPE: 8} (Echo Request) sarebbe necessario catturare i pacchetti in uscita,
ma con \texttt{AF\_INET + SOCK\_RAW} il kernel consegna solo i pacchetti in arrivo
alla macchina locale.

\end{document}
